\section{THUẬT TOÁN K-NN}

\subsection{Khái niệm}
- Là 1 thuật toán học máy có giám sát thường được sử dụng để phân loại.\\
- Hoạt động bằng cách tìm các điểm dữ liệu gần nhất “k”(hàng xóm) với 1 đầu vào nhất định và đưa ra dự đoán dựa trên lớp đa số hoặc giá trị trung bình.

\subsection{Các tham số cốt lõi}
Để vận hành K-NN, ta cần quan tâm hai yếu tố chính:
\begin{itemize}
    \item Tham số $k$ :
    \begin{itemize}
        \item Đây là tham số quyết định số lượng điểm hàng xóm tham gia vào việc bỏ phiếu.
        \item Lựa chọn $k$: Nếu $k$ quá nhỏ (ví dụ $k=1$), mô hình rất nhạy cảm với nhiễu. Nếu $k$ quá lớn, ranh giới phân lớp sẽ bị làm mờ, dẫn đến độ chính xác cục bộ giảm.
    \end{itemize}
    \item Độ đo khoảng cách :
    \begin{itemize}
        \item Để biết đâu là "hàng xóm gần nhất", ta cần thước đo. Phổ biến nhất là Khoảng cách Euclid (Euclidean Distance) - đoạn thẳng nối hai điểm trong không gian hình học.
        \item Công thức tổng quát giữa hai điểm $P$ và $Q$ trong không gian $n$ chiều:
        $$d(P,Q) = \sqrt{\sum_{i=1}^{n}(p_i - q_i)^2}$$
        Trong đó: $p_i$ và $q_i$ lần lượt là giá trị tọa độ (thuộc tính) thứ $i$ của hai điểm dữ liệu.
    \end{itemize}
\end{itemize}

\subsection{Quy trình thực hiện thuật toán}
Quy trình hoạt động của K-NN được chia thành 4 bước tuần tự:
\begin{itemize}
    \item \textbf{Bước 1:} Chuẩn bị tập huấn luyện có đặc trưng và nhãn; tách riêng tập kiểm thử.
    \item \textbf{Bước 2:} Tính khoảng cách từ điểm dữ liệu mới đến tất cả các điểm trong tập huấn luyện.
    \item \textbf{Bước 3:} Sắp xếp các điểm theo khoảng cách tăng dần và xác định $k$ điểm gần nhất.
    \item \textbf{Bước 4:} Dự đoán nhãn cho điểm dữ liệu mới dựa trên nhãn của $k$ điểm gần nhất (bỏ phiếu đa số).
\end{itemize}

\subsection{Mã giả thuật toán}
\begin{lstlisting}[caption={Mã giả thuật toán KNN}]
Algorithm KNN_Prediction:
Input:
    Du lieu Training (X_train, y_train)
    Diem du lieu moi can du doan (x_new)
    Tham so k 
Begin:
    1. Khoi tao danh sach luu tru khoang cach (Distances_List)
    2. Lap qua tung diem du lieu (x_i) trong X_train:
       a. Tinh khoang cach d = Euclidean(x_i, x_new)
       b. Them d vao Distances_List
    3. Sap xep Distances_List theo chieu tang dan.
    4. Chon ra k diem co khoang cach nho nhat (Top k neighbors).
    5. Thong ke nhan cua k diem nay.
    6. Return : Nhan co tan suat xuat hien cao nhat.
End
\end{lstlisting}

\subsection{Ví dụ minh họa chi tiết}
\textbf{Training dataset:}
\begin{table}[H]
    \centering
    \begin{tabular}{|c|c|c|c|}
        \hline
        \textbf{ID} & \textbf{Feature 1} & \textbf{Feature 2} & \textbf{Class} \\
        \hline
        1 & 3 & 4 & A \\ \hline
        2 & 1.5 & 5 & A \\ \hline
        3 & 6 & 8 & A \\ \hline
        4 & 2.5 & 5 & A \\ \hline
        5 & 5 & 2.5 & B \\ \hline
        6 & 5 & 2 & B \\ \hline  
        7 & 8 & 3 & B \\ \hline
        8 & 4.5 & 3 & B \\ \hline
    \end{tabular}
\end{table}

\textbf{Testing dataset} (Size: 10 instances, 2 features, 2 classes):
\begin{table}[H]
    \centering
    \begin{tabular}{|c|c|c|c|}
        \hline
        \textbf{ID} & \textbf{Feature 1} & \textbf{Feature 2} & \textbf{Class} \\
        \hline
        9 & 8 & 1.5 & ? \\ \hline
        10 & 4.5 & 7 & ? \\ \hline
    \end{tabular}
\end{table}

\begin{itemize}
    \item Tính toán cho Test point 9 $(8, 1.5)$ với $k=3$: Sau khi tính khoảng cách và xếp hạng (Rank), 3 láng giềng gần nhất là các mẫu có ID 7, 6, 5 (đều thuộc Class B). $\rightarrow$ \textbf{Kết quả:} Test point 9 có class B.
    \item Tính toán cho Test point 10 $(4.5, 7)$ với $k=3$: Sau khi tính khoảng cách và xếp hạng (Rank), 3 láng giềng gần nhất là các mẫu có ID 3, 8, 7 (đều thuộc Class A). $\rightarrow$ Kết quả: Test point 10 có class A.
\end{itemize}

\subsection{Chương trình cho bài toán minh họa}
\lstinputlisting[language=Python,caption={Chương trình Python minh họa thuật toán KNN}]{source_code/knn66.py}

\subsection{Ưu/nhược điểm thuật toán K-NN}
\begin{itemize}
    \item \textbf{Ưu điểm:}
    \begin{itemize}
        \item \textit{Không yêu cầu mô hình phức tạp:} K-NN không yêu cầu xây dựng mô hình trước mà sử dụng dữ liệu trực tiếp để phân loại.
        \item \textit{Đơn giản và dễ hiểu:} Thuật toán rất đơn giản để triển khai và trực quan, đặc biệt khi dữ liệu ít thuộc tính.
        \item \textit{Tận dụng tính cục bộ:} Có khả năng mô tả tốt các đặc điểm cục bộ vì dựa trên những điểm gần nhất.
    \end{itemize}
    \item \textbf{Nhược điểm:}
    \begin{itemize}
        \item \textit{Tốc độ chậm:} Có thể chậm trong giai đoạn phân loại khi dữ liệu lớn vì phải tính khoảng cách đến tất cả các mẫu.
        \item \textit{Dễ bị ảnh hưởng bởi nhiễu:} Các điểm dữ liệu nhiễu hoặc thuộc tính không liên quan làm giảm độ chính xác.
        \item \textit{Yêu cầu bộ nhớ lớn:} Cần lưu giữ toàn bộ tập dữ liệu huấn luyện để phân loại.
    \end{itemize}
\end{itemize}